\chapter*{ABSTRAK}
\addcontentsline{toc}{chapter}{ABSTRAK}

\begin{center}

    \textbf{Pengembangan Sistem Koreksi Tarif Pajak Penghasilan Berbasis AI pada Web Tukar Faktur di PT Paragon Technology and Innovation } \\
    \vspace{0.4cm}
    oleh \\
    \vspace{0.4cm}
    Wisa Ahmaduta Dinutama\\
    18223003\\
\end{center}


\begin{spacing}{1.0}
PT Paragon Technology and Innovation sebagai perusahaan manufaktur dan distribusi produk kecantikan memiliki proses pengadaan barang dan jasa yang kompleks, termasuk pemenuhan kewajiban perpajakan melalui pemotongan Pajak Penghasilan (PPh).
Penentuan tarif PPh pada transaksi pembelian bergantung pada pemetaan sejumlah atribut, salah satunya jenis material.
Variasi jenis material yang sangat banyak dan sulit dipahami oleh pengguna menyebabkan kesalahan pemilihan material saat pembuatan \textit{Purchase Order} (PO), sehingga tarif pajak yang tercatat menjadi tidak sesuai.
Ketidaksesuaian ini berpotensi menimbulkan kerugian finansial berupa selisih pemotongan pajak, biaya koreksi transaksi, penundaan pembayaran, serta risiko sanksi akibat ketidakpatuhan perpajakan.
Perbaikan langsung pada proses pembuatan PO sulit dilakukan karena berada di luar cakupan kerja divisi penulis dan memerlukan perubahan cara kerja lintas unit yang signifikan.
Untuk mengatasi permasalahan tersebut, penulis mengembangkan Sistem Koreksi Tarif Pajak Penghasilan Berbasis Akal Imitasi (AI) pada \textit{Web Tukar Faktur} sebagai alternatif yang dapat mendeteksi dan mengoreksi ketidaksesuaian tarif PPh tanpa mengubah alur pembuatan PO.
Kerja praktik dilaksanakan di PT Paragon Technology and Innovation, salah satu perusahaan terbesar di industri kecantikan nasional dengan tingkat kematangan teknologi yang tergolong matang sekaligus sedang berkembang pesat, sebagai \textit{software engineer intern} dengan pendekatan pengembangan perangkat lunak yang meliputi analisis kebutuhan, perancangan arsitektur sistem, implementasi \textit{backend}, integrasi layanan AI untuk ekstraksi dan analisis dokumen faktur, serta pengujian fungsional.
Sistem yang dikembangkan menerima data faktur beserta konteks transaksi, membandingkan tarif PPh yang tercatat dengan rekomendasi berbasis aturan dan AI, lalu memberikan hasil koreksi yang dapat digunakan pada proses pertukaran faktur.
Hasil kerja praktik menunjukkan bahwa solusi berbasis AI mampu membantu mengurangi risiko kesalahan tarif PPh akibat pemilihan material yang keliru, sekaligus memberikan pengalaman pengembangan sistem informasi terintegrasi di lingkungan perusahaan.

Kata kunci: akal imitasi, \textit{large language model}, pajak penghasilan, \textit{Purchase Order}, \textit{Enterprise Resource Planning}

\end{spacing}

